% -----------------------------------------------------------------------------
% ВНИМАНИЕ! Все приведённые ниже команды описываются
% в соответствующих разделах appendix/styleguide
% -----------------------------------------------------------------------------
\makeatletter


% =============================================================================
% ---> rules (Общие правила оформления)

\newcommand{\td}[1]{ % Обычная заметка (оранжевая)
    \todo[inline, color=todobackground, linecolor=todoborder, bordercolor=todoborder]{#1}%
}

\newcommand{\fx}[1]{ % Важная заметка (красная)
    \todo[inline, color=xred!20, linecolor=xred, bordercolor=xred]{#1}%
}


% =============================================================================
% ---> Акценты и штрихи

% ПУСТО


% =============================================================================
% ---> Арифметические операции

\AtBeginDocument{
  \renewcommand{\le}{\leqslant}
  \renewcommand{\ge}{\geqslant}
}

\newcommand{\mul}{\mathbin{\raisebox{0.15ex}{\scalebox{0.6}{$\bullet$}}}}

\newcommand{\inc}{\mathord{{+}\mkern-1.5mu{+}}}

\newcommand{\divby}{\mathrel{\vdots}}
\newcommand{\ndivby}{\mathrel{\not\vdots}}

\let\oldsqrt\sqrt
\renewcommand{\sqrt}[2][]{%
  \ifstrempty{#1}{%
    \oldsqrt{#2}%
  }{%
    \oldsqrt[\leftroot{2}\uproot{2}#1]{#2}%
  }%
}

\newcommand{\is}{\coloneqq}
\newcommand{\isr}{\eqqcolon}


% =============================================================================
% ---> Блоки

\newcommand{\noteicon}{\textcolor{xgray}{\faInfoCircle}\;\,}
\newlength{\noteiconwidth}
\newlength{\styletablew}

\newenvironment{styletable}{%
    \setlength{\styletablew}{\dimexpr\textwidth-6\tabcolsep-4\arrayrulewidth\relax}%
  \begin{longtable}{%
    |>{\raggedright\arraybackslash\hyphenpenalty=10000\exhyphenpenalty=10000}p{0.25\styletablew}%
    |p{0.3\styletablew}%
    |p{0.45\styletablew}|}
    \hline
    \textbf{Описание} & \textbf{Пример} & \textbf{Команда} \\
    \hline
    \endfirsthead
    \hline
    \textbf{Описание} & \textbf{Пример} & \textbf{Команда} \\
    \hline
    \endhead
    \hline
    \multicolumn{3}{r}{\footnotesize\itshape\color{xgray}продолжение на следующей странице}\\
    \endfoot
    \hline
    \endlastfoot
}{%
    \end{longtable}%
}

\newenvironment{styletableII}{%
    \setlength{\styletablew}{\dimexpr\textwidth-4\tabcolsep-3\arrayrulewidth\relax}%
    \begin{longtable}{|p{0.3\styletablew}|p{0.7\styletablew}|}
    \hline
    \textbf{Элемент} & \textbf{Синтаксис} \\
    \hline
    \endfirsthead
    \hline
    \textbf{Элемент} & \textbf{Синтаксис} \\
    \hline
    \endhead
    \hline
    \multicolumn{2}{r}{\footnotesize\itshape\color{xgray}продолжение на следующей странице}\\
    \endfoot
    \hline
    \endlastfoot
}{%
    \end{longtable}%
}

% См. также файл block.tex


% =============================================================================
% ---> Выделение

%\newcommand{\foc}[1]{%
%  \textcolor{xstone!70!black}{\bfseries\underLine{#1}}%
%}
\newcommand{\foc}[1]{%
  \underLine{{\bfseries\color{xstone!70!black}#1}}%
}

\newcommand{\emp}[1]{%
  \textbf{\textcolor{xred!70!black}{#1}}%
}

\newcommand{\stylenote}[1]{%
  \par\vspace{0.7em}%
  \noindent
  \begin{tcolorbox}[
    enhanced,
    breakable,
    colback=xgray!6,
    colframe=xgray!40,
    boxrule=0.4pt,
    arc=1.5mm,
    left=2.5mm, right=2.5mm, top=2mm, bottom=2mm,
    before skip=4mm plus 3mm minus 2mm,
    after skip=5mm plus 3mm minus 2mm,
    leftrule=2.5pt,
    rightrule=0.4pt,
    toprule=0.4pt,
    bottomrule=0.4pt,
    parbox=false,
    before upper app={%
      \small
      \settowidth{\noteiconwidth}{\noteicon}%
      \setlength{\parindent}{\noteiconwidth}%
      \setlength{\parskip}{0.5em plus 0.15em minus 0.05em}%
    },
  ]
    \noindent\noteicon #1%
  \end{tcolorbox}%
}


% =============================================================================
% ---> Интегралы

% ПУСТО


% =============================================================================
% ---> Логика и кванторы

\newcommand{\im}{\mathrel{\DOTSB\Rightarrow}}
\newcommand{\eq}{\mathrel{\DOTSB\Leftrightarrow}}

\newcommand{\existso}{\exists!\,}

\newcommand{\qs}{\mathrel{:}}

\AtBeginDocument{%
  \providecommand{\eqdef}{}%
  \renewcommand{\eqdef}{\mathrel{\vcentcolon\Leftrightarrow}}%
}
\newcommand{\eqdefr}{\mathrel{\Leftrightarrow\vcentcolon}}


% =============================================================================
% ---> Множества

\newcommand{\set}[1]{\left\{ \, #1 \, \right\}}
\newcommand{\setc}[2]{\left\{ \, #1 \;\middle|\; #2 \, \right\}}

\newcommand{\card}[1]{\left| \, #1 \, \right|}

\AtBeginDocument{\renewcommand{\emptyset}{\varnothing}}

\newcommand{\seg}[2]{\left[#1, \, #2\right]}
\newcommand{\sego}[2]{\left(#1, \, #2\right)}
\newcommand{\segor}[2]{\left[#1, \, #2\right)}
\newcommand{\segol}[2]{\left(#1, \, #2\right]}

\newcommand{\tup}[1]{\left(#1\right)}


% =============================================================================
% ---> Отображения

% ПУСТО


% =============================================================================
% ---> Пределы

\newcommand{\nh}[2][\e]{U_{#1}(#2)}             % Окрестность
\newcommand{\nhp}[2][\e]{\mathring{U}_{#1}(#2)} % Проколотая окрестность

\newcommand{\appr}[1]{\xrightarrow[#1]{}}       % Предел со стрелкой

\newcommand{\uplim}{\varlimsup\limits}
\newcommand{\lowlim}{\varliminf\limits}

\DeclareMathOperator{\olandau}{o}
\newcommand{\so}[1]{\olandau\mathopen{}\left(#1\right)\mathclose{}}
\DeclareMathOperator{\Olandau}{O}
\newcommand{\bo}[1]{\Olandau\mathopen{}\left(#1\right)\mathclose{}}


% =============================================================================
% ---> Производные

\newcommand{\dif}{\mathop{}\!\mathrm{d}}

\newcommand{\od}[3][1]{%
  \ifstrequal{#1}{1}%
    {\frac{\dif #2}{\dif #3}}%
    {\frac{\dif^{#1} #2}{\dif #3^{#1}}}%
}

\newcommand{\pd}[3][1]{%
  \ifstrequal{#1}{1}%
    {\frac{\partial #2}{\partial #3}}%
    {\frac{\partial^{#1} #2}{\partial #3^{#1}}}%
}


% =============================================================================
% ---> Ряды

% ПУСТО


% =============================================================================
% ---> Символы

\newcommand{\e}{\varepsilon}

\renewcommand{\d}{\delta}
\newcommand{\D}{\Delta}

\newcommand{\rom}[1]{\text{\MakeUppercase{\expandafter\romannumeral\number#1\relax}}}


% =============================================================================
% ---> Сокращения

\newcommand{\tit}{тогда и только тогда, когда\xspace}


% =============================================================================
% ---> Числовые множества

\newcommand{\N}{\mathbb{N}}
\newcommand{\Nz}{\mathbb{N}_0}
\newcommand{\Z}{\mathbb{Z}}
\newcommand{\Q}{\mathbb{Q}}
\newcommand{\I}{\mathbb{I}}
\newcommand{\R}{\mathbb{R}}
\newcommand{\Rext}{\overline{\mathbb{R}}}
\newcommand{\C}{\mathbb{C}}


% =============================================================================
% ---> Числовые характеристики

\newcommand{\abs}[1]{\left|{#1}\right|}

\newcommand{\norm}[1]{\lVert{#1}\rVert}

\newcommand{\floor}[1]{\left\lfloor #1 \right\rfloor}
\newcommand{\ceil}[1]{\left\lceil #1 \right\rceil}
\newcommand{\fpart}[1]{\left\{ #1 \right\}}

\DeclareMathOperator{\sgn}{sgn}

\DeclareMathOperator{\Cre}{Re}
\DeclareMathOperator{\Cim}{Im}


% =============================================================================
% ---> Экстремумы и грани

% ПУСТО


% =============================================================================
% ---> Элементарные функции

\DeclareMathOperator{\tg}{tg}
\DeclareMathOperator{\ctg}{ctg}
\DeclareMathOperator{\arctg}{arctg}
\DeclareMathOperator{\arcctg}{arcctg}


% =============================================================================
\makeatother